\chapter{Conclusion and Future Work}
\label{chap:conclusion}

\section{Answers to the research questions}

\subsection{RQ1: leakage-controlled data handoff}

The project produces chronological, non-overlapping model-ready splits and
persists the transformation artifacts needed by downstream training. The saved
verifier reports 82 passed checks and no failures. The only identified lifecycle
defect is the stale terminal status in the top-level manifest.

\subsection{RQ2: model comparison}

LightGBM has the highest saved validation PR-AUC in both V1 and V2 comparisons.
The V2 final artifact reports a holdout ROC-AUC of 0.8796 and PR-AUC of 0.4582,
improving on V1. Because V1 and V2 differ in feature contract and because
operating-point metrics are absent, V2 should be described as the offline
champion rather than proof of production benefit.

\subsection{RQ3: Spark-free serving and explanations}

The model package serializes categorical mappings and feature order so FastAPI
can score without Spark. Supported tree artifacts return local SHAP
contributions. This design successfully separates batch feature preparation
from request-time inference, subject to the availability of the same engineered
features in a real online environment.

\subsection{RQ4: production-readiness gaps}

The principal gaps are:

\begin{itemize}
  \item incomplete experiment/model lineage;
  \item missing calibration and operating-policy evidence;
  \item fail-open heuristic fallback;
  \item no authentication, durable jobs, or migrations;
  \item no drift/latency/service monitoring;
  \item manual client contract synchronization;
  \item uneven CI and frontend test coverage.
\end{itemize}

\section{Technical contributions}

The system demonstrates a complete flow from multi-gigabyte source data to a
human review decision:

\begin{enumerate}
  \item Spark performs typed ingestion, audits, EDA, chronological splitting,
  and leakage-controlled feature engineering.
  \item Explicit contracts make the data--model handoff inspectable.
  \item Multiple classifiers are compared under an imbalanced metric.
  \item The selected estimator is packaged with categorical mappings and SHAP.
  \item FastAPI, SQLAlchemy, and React implement an explainable review workflow.
  \item Docker Compose provides local application and training profiles.
\end{enumerate}

\section{Immediate future work}

The first release-hardening milestone should not train a new model. It should
make V2 independently reproducible and observable:

\begin{enumerate}
  \item complete dependency locking and a frozen container smoke;
  \item record model checksum, image digest, data/schema versions, and MLflow
  run;
  \item calculate operating-point and calibration evidence;
  \item make fallback an explicit demo mode and enforce readiness;
  \item synchronize/generated API client types and expand CI.
\end{enumerate}

\section{Model-development future work}

After reproducibility:

\begin{itemize}
  \item perform temporal cross-validation;
  \item separate model selection, calibration, threshold selection, and holdout;
  \item run feature-group ablations;
  \item evaluate paired V1/V2 errors by segment;
  \item measure inference latency and resource consumption;
  \item monitor delayed-label performance and drift.
\end{itemize}

\section{Platform future work}

A deployable platform requires:

\begin{itemize}
  \item authentication, authorization, secrets, and immutable audit events;
  \item Alembic migrations and durable background queues;
  \item a serving-only model dependency profile;
  \item production metrics, logs, tracing, readiness, and alerts;
  \item bulk re-scoring or version-isolated reporting;
  \item retention, privacy, and incident-response policies.
\end{itemize}

\section{Final statement}

The project is a coherent and well-integrated academic fraud risk scoring
demonstration. Its strongest contribution is not a single metric; it is the
connection of data contracts, a trained model, explanations, API behavior,
persistence, and human review. The whole-project audit defines the additional
evidence and controls required to turn that demonstration into a trustworthy
deployment candidate.
